%==============================================================================
% Dàn ý chi tiết - Chapter 4: Thực nghiệm và đánh giá
%==============================================================================

\chapter{Thực nghiệm và đánh giá}

%-----------------------------------------------------------------------------
\section{Thiết lập thực nghiệm}
%-----------------------------------------------------------------------------

\subsection{Các mô hình sử dụng}
\begin{itemize}
    \item Danh sách các LLM được thử nghiệm:
    \begin{itemize}
        \item GPT-4 (OpenAI)
        \item GPT-3.5 Turbo
        \item Claude-3 (Opus/Sonnet/Haiku)
        \item LLaMA-3 (8B, 70B)
        \item Mistral/Mixtral
    \end{itemize}
    \begin{itemize}
        \item Version cụ thể
    \end{itemize}
    \begin{itemize}
        \item Nguồn truy cập (API, local deployment)
    \end{itemize}
    \begin{itemize}
        \item So sánh các model về:
        \begin{itemize}
            \item Pricing
            \item Capability
            \item Availability
        \end{itemize}
    \end{itemize}
    \item Các mô hình OCR:
    \begin{itemize}
        \item EasyOCR
        \item PaddleOCR
        \item Tesseract
    \end{itemize}
    \begin{itemize}
        \item GPT-4V / Claude Vision
    \end{itemize}
\end{itemize}

\subsection{Cấu hình thực nghiệm}
\begin{itemize}
    \item Hardware:
    \begin{itemize}
        \item Server specifications (nếu có local deployment)
    \end{itemize}
    \item Software environment:
    \begin{itemize}
        \item Python version
        \item Key libraries (LangGraph, LangChain, etc.)
    \end{itemize}
    \item Hyperparameters:
    \begin{itemize}
        \item Temperature settings
        \item Max tokens
        \item Top-p
        \item Number of generations (N)
    \end{itemize}
    \item Dataset splits:
    \begin{itemize}
        \item Số lượng samples cho test
    \end{itemize}
    \item Experimental settings:
    \begin{itemize}
        \item Số lượng câu hỏi sinh ra mỗi văn bản
    \end{itemize}
    \item Cost analysis
\end{itemize}

%-----------------------------------------------------------------------------
\section{Kết quả sinh câu hỏi}
%-----------------------------------------------------------------------------

\subsection{Kết quả đánh giá vi mô (Micro Evaluation)}
\begin{itemize}
    \item Solvability Results:
    \begin{itemize}
        \item Tỷ lệ câu hỏi có thể giải được theo từng cấp độ
        \begin{itemize}
            \item Level 1: X\%
            \item Level 2: Y\%
            \item Level 3: Z\%
        \end{itemize}
        \item So sánh giữa các model
    \end{itemize}
    \item Distractor Quality Results:
    \begin{itemize}
        \item Điểm trung bình chất lượng distractors
        \begin{itemize}
            \item Theo thang điểm 1-5
        \end{itemize}
        \item Phân tích theo cấp độ
    \end{itemize}
    \item Alignment Results:
    \begin{itemize}
        \item Tỷ lệ câu hỏi đúng cấp độ Bloom
        \begin{itemize}
            \item Level 1 alignment: X\%
            \item Level 2 alignment: Y\%
            \item Level 3 alignment: Z\%
        \end{itemize}
        \item Confusion matrix (nếu có misclassification)
    \end{itemize}
    \item Nhận xét và phân tích
    \begin{itemize}
        \item Cấp độ nào có kết quả tốt/xấu nhất
        \item Nguyên nhân
    \end{itemize}
\end{itemize}

\subsection{Kết quả đánh giá vĩ mô -- So sánh giữa các mô hình (Macro Evaluation)}
\begin{itemize}
    \item Overall Quality Scores:
    \begin{itemize}
        \item Bảng so sánh điểm trung bình (1-5)
        \begin{itemize}
            \item GPT-4 vs Claude-3 vs LLaMA-3...
        \end{itemize}
    \end{itemize}
    \item Chi tiết từng tiêu chí:
    \begin{itemize}
        \item Grammar and fluency
        \item Clarity
        \item Difficulty appropriateness
        \item Relevance to passage
    \end{itemize}
    \item Statistical analysis:
    \begin{itemize}
        \item Standard deviation
        \item Significant differences (p-value)
    \end{itemize}
    \item Best performing model(s)
    \item Phân tích chi tiết từng cấp độ Bloom
\end{itemize}

\subsection{So sánh phương pháp đơn prompt và đa tác tử}
\begin{itemize}
    \item So sánh Micro Evaluation:
    \begin{itemize}
        \item Solvability
        \item Distractor Quality
        \item Alignment
    \end{itemize}
    \item So sánh Macro Evaluation:
    \begin{itemize}
        \item Overall quality scores
    \end{itemize}
    \item Format evaluation:
    \begin{itemize}
        \item Error rates
    \end{itemize}
    \item Cost comparison:
    \begin{itemize}
        \item API calls
        \item Total cost
        \item Latency
    \end{itemize}
    \item Nhận xét:
    \begin{itemize}
        \item Khi nào nên dùng Single-Prompt
        \item Khi nào nên dùng Multi-Agent
    \end{itemize}
    \item Trade-offs analysis
\end{itemize}

%-----------------------------------------------------------------------------
\section{Kết quả đánh giá định dạng đầu ra}
%-----------------------------------------------------------------------------

\subsection{Tỷ lệ lỗi cú pháp (Syntax Error Rate)}
\begin{itemize}
    \item Định nghĩa lỗi cú pháp:
    \begin{itemize}
        \item Missing question mark
        \item Invalid option format
        \item Incorrect answer indicator
    \end{itemize}
    \item Kết quả:
    \begin{itemize}
        \item Tỷ lệ lỗi theo model
        \item Tỷ lệ lỗi theo cấp độ
    \end{itemize}
    \item Phân tích nguyên nhân lỗi
    \item So sánh Single-Prompt vs Multi-Agent
\end{itemize}

\subsection{Tỷ lệ sinh thừa và sinh thiếu (Over/Under-generation)}
\begin{itemize}
    \item Định nghĩa:
    \begin{itemize}
        \item Over-generation: Sinh nhiều hơn yêu cầu
        \item Under-generation: Sinh ít hơn yêu cầu
    \end{itemize}
    \item Kết quả:
    \begin{itemize}
        \item Thống kê phân phối số câu hỏi sinh ra
    \end{itemize}
    \item Ảnh hưởng đến chất lượng
    \item Cách xử lý trong Merge Agent
\end{itemize}

%-----------------------------------------------------------------------------
\section{Kết quả OCR}
%-----------------------------------------------------------------------------

\subsection{Kết quả trên ảnh sạch}
\begin{itemize}
    \item CER (Character Error Rate):
    \begin{itemize}
        \item EasyOCR: X\%
        \item PaddleOCR: Y\%
        \item Tesseract: Z\%
        \item GPT-4V: W\%
    \end{itemize}
    \item WER (Word Error Rate):
    \begin{itemize}
        \item So sánh tương tự
    \end{itemize}
    \item Nhận xét:
    \begin{itemize}
        \item Model nào tốt nhất cho clean images
    \end{itemize}
    \item Ví dụ errors (nếu có)
\end{itemize}

\subsection{Kết quả trên ảnh tăng cường}
\begin{itemize}
    \item CER/WER trên augmented images:
    \begin{itemize}
        \item So sánh giữa các phương pháp augmentation
        \begin{itemize}
            \item Noise
            \item Blur
            \item Rotation
            \item Combined effects
        \end{itemize}
    \end{itemize}
    \item Impact of augmentation level
    \item Comparison: Traditional OCR vs VLM
    \begin{itemize}
        \item VLM có robust hơn không?
    \end{itemize}
    \item Nhận xét về robustness
    \begin{itemize}
        \item Các loại augmentation ảnh hưởng nhiều nhất
    \end{itemize}
    \item Trade-offs: Accuracy vs Speed vs Cost
\end{itemize}

%-----------------------------------------------------------------------------
\section{Thảo luận}
%-----------------------------------------------------------------------------
\begin{itemize}
    \item Tổng hợp các phát hiện chính
    \begin{itemize}
        \item Đánh giá tổng quan về phương pháp đề xuất
    \end{itemize}
    \item Phân tích sâu hơn:
    \begin{itemize}
        \item Tại sao LLM sinh tốt/các câu hỏi cấp độ nào
        \item Ảnh hưởng của prompt design
        \item Vai trò của Multi-Agent architecture
    \end{itemize}
    \item So sánh với các nghiên cứu liên quan:
    \begin{itemize}
        \item Phương pháp nào tốt hơn?
        \item Điểm mạnh và hạn chế
    \end{itemize}
    \item Các yếu tố ảnh hưởng đến kết quả:
    \begin{itemize}
        \item Chất lượng văn bản đầu vào
        \item Độ dài văn bản
        \subject Complexity of content
    \end{itemize}
    \item Practical implications:
    \begin{itemize}
        \item Recommendations cho practitioners
    \end{itemize}
\end{itemize}

%-----------------------------------------------------------------------------
\section{Hạn chế và hướng phát triển}
%-----------------------------------------------------------------------------

\subsection{Hạn chế của nghiên cứu}
\begin{itemize}
    \item Hạn chế về dữ liệu:
    \begin{itemize}
        \item Chỉ sử dụng RACE dataset
        \item Limited language coverage (tiếng Anh)
    \end{itemize}
    \item Hạn chế về phương pháp:
    \begin{itemize}
        \item Chỉ 3 cấp độ Bloom
        \item Prompt engineering có thể chưa tối ưu
    \end{itemize}
    \item Hạn chế về đánh giá:
    \begin{itemize}
        \item LLM-as-a-Judge có potential bias
        \item Không có human evaluation đầy đủ
    \end{itemize}
    \item Hạn chế về tài nguyên:
    \begin{itemize}
        \item Không test tất cả models
        \item Limited scalability testing
    \end{itemize}
    \item Các giả định và simplifications
\end{itemize}

\subsection{Hướng phát triển tương lai}
\begin{itemize}
    \item Mở rộng phạm vi:
    \begin{itemize}
        \item Thêm nhiều cấp độ Bloom
        \item Đa ngôn ngữ
        \item Nhiều loại câu hỏi (True/False, Fill-in-blank)
    \end{itemize}
    \item Cải thiện phương pháp:
    \begin{itemize}
        \item Fine-tuning models cho question generation
        \item Self-consistency và ensemble methods
        \item Feedback loop cho iterative improvement
    \end{itemize}
    \item Cải thiện đánh giá:
    \begin{itemize}
        \item Human evaluation với larger scale
        \item Cross-annotator agreement
        \item Student performance evaluation (A/B testing)
    \end{itemize}
    \item Ứng dụng thực tế:
    \begin{itemize}
        \item Integration với LMS
        \item Real-time generation
        \item Adaptive testing
    \end{itemize}
    \item Nghiên cứu liên quan:
    \begin{itemize}
        \item Knowledge tracking
        \item Item response theory integration
    \end{itemize}
\end{itemize}
